\documentclass[11pt]{article}
\usepackage[a4paper,margin=27mm]{geometry}
\usepackage{amsmath,amssymb,amsthm,mathtools}
\usepackage{booktabs,array}
\usepackage[T1]{fontenc}
\usepackage[utf8]{inputenc}
\usepackage{lmodern}
\usepackage{microtype}
\usepackage{xcolor}
\usepackage{hyperref}
\hypersetup{colorlinks=true,linkcolor=blue!55!black,citecolor=blue!55!black,urlcolor=blue!55!black}
\setlength{\parindent}{0pt}
\setlength{\parskip}{5pt plus 1pt}
\newtheorem{theorem}{Theorem}[section]
\newtheorem{lemma}[theorem]{Lemma}
\newtheorem{proposition}[theorem]{Proposition}
\newtheorem{corollary}[theorem]{Corollary}
\theoremstyle{remark}
\newtheorem{remark}[theorem]{Remark}
\newcommand{\per}{\operatorname{per}}
\newcommand{\one}{\mathbf 1}
\newcommand{\RR}{\mathbb R}
\newcommand{\QQ}{\mathbb Q}

\title{Dittert's conjecture in dimension 10\\
\large An exact computer-assisted boundary-exclusion proof}
\author{Pedro Paulo Marques do Nascimento}
\date{23 July 2026}

\begin{document}
\maketitle

\begin{center}
\fbox{\parbox{0.90\textwidth}{\textbf{Status.} This is an exact, reproducible
working proof. It has not yet undergone independent peer review. The finite
dimension-specific calculations are checked by two separately implemented
standard-library verifiers.}}

\small Licensed under \href{https://creativecommons.org/licenses/by/4.0/}{CC BY 4.0}.
\end{center}

\begin{abstract}
For a nonnegative $n\times n$ matrix $A$ whose entries sum to $n$, let $r_i,c_j$
be its row and column sums and put
\[
 \Phi(A)=\prod_i r_i+\prod_j c_j-\per A.
\]
We prove Dittert's conjecture in dimension $n=10$.  As in the proofs for
$n=11,12,13$, a maximal Li scaling reduces the boundary problem to marginal
and proper tight cuts.  Proper cuts are excluded by exact staircase-face
bounds.  For a marginal cut we add the stationarity equation at its minimum
row or column.  It forces a line-sum deviation on the opposite side, allowing
the shared row--column product deficit to exclude the last case.  The two
univariate inequalities are certified exactly, first by rational Sturm
sequences and independently by exact Ryser interpolation and Bernstein
positivity.
\end{abstract}

\section{Statement and published inputs}

Let
\[
 K_n=\left\{A=(a_{ij})\in\RR_{\ge0}^{n\times n}:
       \sum_{i,j}a_{ij}=n\right\},
 \qquad
 \gamma_n=\frac{n!}{n^n}.
\]
For $A\in K_n$, write
\[
 R=\prod_{i=1}^n r_i,\qquad
 C=\prod_{j=1}^n c_j,\qquad
 P=\per A.
\]
Thus $\Phi(A)=R+C-P$ and
\[
 \Phi(J_n/n)=2-\gamma_n.
\]

\begin{theorem}\label{thm:main}
For every $A\in K_{10}$,
\[
 \Phi(A)\le 2-\frac{10!}{10^{10}},
\]
with equality if and only if $A=J_{10}/10$.
\end{theorem}

We use the following published results.
\begin{enumerate}
\item Hwang proved that the only possible positive $\Phi$-maximizer in $K_n$
is $J_n/n$ \cite{Hwang1986}.
\item Cheon and Wanless proved that a partly decomposable matrix cannot
maximize $\Phi$, and that a maximizer cannot have its complete zero set equal,
up to row and column permutations, to one proper rectangular block
\cite[Theorems 3.2 and 3.7]{CheonWanless2012}.
\item Li's cut characterization says that a nonnegative $n\times n$ matrix
$S$ is doubly superstochastic if and only if
\[
 \sum_{i\in I,j\in J}s_{ij}\ge |I|+|J|-n
\]
for all $I,J\subseteq[n]$ \cite{Li1990}; see also
\cite[Lemma 2.2]{CheonWanless2012}.
\item On a staircase face of the Birkhoff polytope, the barycenter minimizes
the permanent \cite{Hwang1985}.
\item Minc's averaging theorem for permanent minimizers with repeated support
rows and columns gives the block form used below \cite[Theorem 1]{Minc1984};
see also \cite[Section 2]{PulaSongWanless2011}.
\item The van der Waerden permanent theorem gives $\per B\ge\gamma_n$ for
every doubly stochastic $B$, with equality only at $J_n/n$
\cite{Egorychev1981,Falikman1981}.
\end{enumerate}
The stationarity calculation that is new to the marginal-cut argument is
proved directly in Section~\ref{sec:marginal}; it is not an additional
external assumption.

\section{Boundary structure and the shared deficit}

Because $K_n$ is compact, $\Phi$ has a global maximizer.  By Hwang's
positive-support theorem, a nonuniform maximizer must lie on the boundary.

\begin{lemma}[Two independent zeroes]\label{lem:independent-zeros}
Every boundary $\Phi$-maximizer has two zero entries in distinct rows and
distinct columns.
\end{lemma}
\begin{proof}
Suppose a nonempty zero set contains no such pair.  Fix a zero at $(i,j)$.
Every other zero shares row $i$ or column $j$.  If there were both a zero
$(i,j')$, $j'\ne j$, and a zero $(i',j)$, $i'\ne i$, those two zeroes would be
in distinct rows and columns.  Hence all zeroes lie in one row or all lie in
one column.

If that row or column is entirely zero, the matrix is partly decomposable.
Otherwise its complete zero set is one proper rectangular block.  Both cases
are excluded for a maximizer by Cheon--Wanless.
\end{proof}

Assume for contradiction that a boundary maximizer $A$ satisfies
\[
 \Phi(A)\ge 2-\gamma_n.
\]
The row sums and column sums each have arithmetic mean $1$, so AM--GM gives
$R,C\le1$.  Consequently $P\le\gamma_n$.  Write
\[
 P=\gamma_n-\delta,qquad 0\le\delta\le\gamma_n,
\]
and put $R=1-\rho$, $C=1-\sigma$.  The assumed lower bound on $\Phi$ gives
\begin{equation}\label{eq:shared-deficit}
 \rho+\sigma\le\delta.
\end{equation}
In particular,
\begin{equation}\label{eq:product-lower}
 R\ge1-\delta,\qquad C\ge1-\delta,
 \qquad RC=(1-\rho)(1-\sigma)\ge1-\delta.
\end{equation}

\section{The maximal Li scaling}

For $I,J\subseteq[n]$, let $s(A[I,J])$ be the sum of the indicated
submatrix, and define
\begin{equation}\label{eq:lambda}
 \lambda=
 \min_{\substack{I,J\subseteq[n]\\k:=|I|+|J|-n>0}}
 \frac{s(A[I,J])}{k}.
\end{equation}
A maximizer is fully indecomposable by the Cheon--Wanless exclusion.  Hence
every numerator in \eqref{eq:lambda} is positive, so $\lambda>0$.  Taking
$I=J=[n]$ shows $\lambda\le1$.

Li's theorem says that $A/\lambda$ is doubly superstochastic.  Thus there is
a doubly stochastic matrix $B$ such that
\begin{equation}\label{eq:domination}
 \lambda B\le A\quad\hbox{entrywise},
 \qquad
 P=\per A\ge\lambda^n\per B.
\end{equation}
Choose $I,J$ attaining \eqref{eq:lambda}, and set
\[
 u=n-|I|,\qquad v=n-|J|,\qquad k=n-u-v>0.
\]
There are three cases: the whole cut $(u,v)=(0,0)$, a marginal cut where
exactly one of $u,v$ is zero, and a proper cut where $u,v\ge1$.

If $u=v=0$, then $\lambda=1$.  Since $B\le A$ and both matrices have total
sum $n$, $A=B$ is doubly stochastic.  The van der Waerden theorem and
$P\le\gamma_n$ then force $A=J_n/n$, contrary to the boundary assumption.

\section{The two-independent-zero permanent gap}

Let $\Omega_n^{(2)}$ be the face of the Birkhoff polytope specified by
$b_{12}=b_{21}=0$, and put $m=n-2$.  Minc's averaging theorem permits a
permanent minimizer on this face to be chosen in the form
\begin{equation}\label{eq:Y}
 Y=
 \begin{pmatrix}
 x_1&0&a_1\one^{\mathsf T}\\
 0&x_2&a_2\one^{\mathsf T}\\
 a_1\one&a_2\one&xJ_m
 \end{pmatrix},
 \qquad
 x_i=1-ma_i,
 \qquad
 a_1+a_2+mx=1.
\end{equation}
Here $0\le a_i\le1/m$.  We include the equalization step, since only the
averaging theorem is imported.

\begin{proposition}[One-variable reduction]\label{prop:face-reduction}
For $m\ge4$, the minimum permanent in \eqref{eq:Y} occurs with $a_1=a_2$.
Consequently the minimum of $\per B$ on $\Omega_n^{(2)}$ is the minimum, on
\[
 \frac{m-2}{m^2}\le x\le\frac1m,
\]
of
\begin{equation}\label{eq:Pn}
 P_n(x)=m!x^{m-2}
 \left(x^2a^2+2mxab^2+m(m-1)b^4\right),
\end{equation}
where
\begin{equation}\label{eq:ab}
 a=\frac{2-m+m^2x}{2},\qquad b=\frac{1-mx}{2}.
\end{equation}
\end{proposition}
\begin{proof}
Put $s=a_1+a_2$, $p=a_1a_2$, and $x=(1-s)/m$.  Counting permanent terms in
\eqref{eq:Y} gives, apart from the positive factor $m!x^{m-2}$,
\begin{equation}\label{eq:Esp}
 E_m(s,p)=x^2(1-ms+m^2p)
 +mx\bigl(s^2-(2+ms)p\bigr)+m(m-1)p^2.
\end{equation}
For fixed $s$,
\begin{equation}\label{eq:Ep}
 \frac{\partial E_m}{\partial p}
 =2m(m-1)p+(m+1)s^2-ms-1.
\end{equation}
Since $p\le s^2/4$ and $0\le s\le2/m$, the right side is at most
\[
 D_m(s)=\frac{(m^2+m+2)s^2-2ms-2}{2}.
\]
This quadratic is convex, while
\[
 D_m(0)=-1,
 \qquad
 D_m(2/m)=-1+\frac2m+\frac4{m^2}<0
 \quad(m\ge4).
\]
Thus \eqref{eq:Ep} is strictly negative.  For fixed $s$, the permanent is
minimized by maximizing $p$, namely by $p=s^2/4$ and $a_1=a_2=b$.  Solving
the stochastic constraints yields \eqref{eq:ab} and the interval for $x$.
A second count gives \eqref{eq:Pn}.
\end{proof}

For $n=10$, set
\begin{equation}\label{eq:L}
 L_{10}=\frac{36719}{10^8}.
\end{equation}

\begin{proposition}[Exact two-zero bound]\label{prop:two-zero-bound}
Every $B\in\Omega_{10}^{(2)}$ satisfies
\[
 \per B>L_{10}>\gamma_{10}.
\]
\end{proposition}
\begin{proof}
Here $m=8$ and the complete interval is $[3/32,1/8]$.  The primary verifier
expands $P_{10}(x)-L_{10}$ over $\QQ$, verifies positive endpoint values, and
uses an exact rational Sturm sequence to prove that it has no root in the
interval.

The independent verifier does not use that expansion.  It evaluates the
permanent of the matrix \eqref{eq:ab} by the literal exact Ryser formula at
eleven rational points.  Degree at most ten makes these values determine the
polynomial.  On each of 256 equal rational subintervals it converts the
polynomial to the degree-ten Bernstein basis; every one of the resulting
coefficients is strictly positive.  This independently proves
$P_{10}(x)-L_{10}>0$ on the entire interval.  Finally,
\[
 L_{10}-\gamma_{10}=\frac{431}{10^8}>0.
\]
\end{proof}

\section{Marginal tight cuts}\label{sec:marginal}

We first record the elementary stationarity equation used in this case.
Since a maximizer is fully indecomposable, all row and column sums are
positive.

\begin{lemma}[Support stationarity]\label{lem:stationarity}
At a global maximizer, for every positive entry $a_{ij}$,
\begin{equation}\label{eq:stationarity}
 \frac{R}{r_i}+\frac{C}{c_j}-\per A(i\mid j)=\Phi(A),
\end{equation}
where $A(i\mid j)$ is obtained by deleting row $i$ and column $j$.
\end{lemma}
\begin{proof}
The left side of \eqref{eq:stationarity} is the partial derivative
$\partial\Phi/\partial a_{ij}$.  At a maximum on the nonnegative simplex,
moving a small amount of mass between two positive entries shows that all
partial derivatives on the support have a common value $\mu$; moving mass
from a positive entry to a zero entry shows only the corresponding one-sided
inequality at a zero.  Euler's identity applies because $\Phi$ is homogeneous
of degree $n$:
\[
 n\Phi(A)=\sum_{i,j}a_{ij}\frac{\partial\Phi}{\partial a_{ij}}
 =\mu\sum_{i,j}a_{ij}=n\mu.
\]
Thus $\mu=\Phi(A)$.
\end{proof}

Suppose the tight cut is marginal.  The ratio in \eqref{eq:lambda} is then
an average of a nonempty set of row sums or column sums.  Singleton marginal
cuts are also allowed, so
\[
 \lambda=\min_i r_i
 \quad\hbox{or}\quad
 \lambda=\min_j c_j.
\]
By transposition, it is enough to treat a minimum row.  Choose $i$ with
\[
 r_i=\lambda=1-a,
 \qquad 0\le a<1.
\]
Multiply \eqref{eq:stationarity} by $a_{ij}$ and sum over $j$.  Expansion of
the permanent along row $i$ gives
\[
 R+C\sum_j\frac{a_{ij}}{c_j}-P=r_i\Phi(A).
\]
Therefore
\begin{equation}\label{eq:q}
 q_i:=\sum_j\frac{a_{ij}}{c_j}
 =1-(1+\theta)a,
 \qquad
 \theta=\frac{R-P}{C}.
\end{equation}
By \eqref{eq:product-lower} and $P=\gamma_n-\delta$,
\begin{equation}\label{eq:theta}
 \theta\ge R-P\ge1-\gamma_n.
\end{equation}
Let $c_{\max}=\max_jc_j$.  Since $q_i>0$, equations
\eqref{eq:q}--\eqref{eq:theta} imply
\[
 \frac{1-a}{c_{\max}}le q_i
 \le1-(2-\gamma_n)a.
\]
The last denominator is therefore positive, and
\begin{equation}\label{eq:h}
 c_{\max}\ge1+h(a),
 \qquad
 h(a)=\frac{(1-\gamma_n)a}{1-(2-\gamma_n)a}.
\end{equation}

AM--GM on the other nine row sums, followed by AM--GM on the other nine
column sums, now gives
\begin{align}
 R&\le f(a):=(1-a)\left(1+\frac a9\right)^9,
 \label{eq:f}\\
 C&\le g(h(a)):=(1+h(a))\left(1-\frac{h(a)}9\right)^9.
 \label{eq:g}
\end{align}
For the second inequality, the product at fixed largest column sum is bounded
by equalizing the other columns; the resulting expression decreases once the
largest sum is at least $1$.

Define
\begin{equation}\label{eq:D}
 D(a)=2-f(a)-g(h(a)).
\end{equation}
The shared deficit \eqref{eq:shared-deficit} and
\eqref{eq:f}--\eqref{eq:g} imply
\begin{equation}\label{eq:delta-D}
 \delta\ge D(a).
\end{equation}
Moreover, $f$ is strictly decreasing for $a>0$, $h$ is strictly increasing,
and $g$ is strictly decreasing for positive arguments.  Hence $D$ is strictly
increasing for $a>0$.  The exact check
\begin{equation}\label{eq:D-endpoint}
 D(1/50)-\gamma_{10}>0
\end{equation}
and $\delta\le\gamma_{10}$ confine every feasible $a$ to
$0\le a<1/50$.

The dominated doubly stochastic matrix $B$ in \eqref{eq:domination} inherits
the two independent zeroes of $A$.  Proposition~\ref{prop:two-zero-bound}
therefore gives
\begin{equation}\label{eq:marginal-permanent}
 P\ge L_{10}(1-a)^{10}.
\end{equation}
On the other hand, \eqref{eq:delta-D} gives
$P=\gamma_{10}-\delta\le\gamma_{10}-D(a)$.  Thus a contradiction follows
once
\begin{equation}\label{eq:H}
 H(a):=L_{10}(1-a)^{10}-\gamma_{10}+D(a)>0
 \qquad(0\le a\le1/50)
\end{equation}
is established.

Put $d(a)=1-(2-\gamma_{10})a$.  Clearing the positive denominator
$9^9d(a)^{10}$ in \eqref{eq:H} gives a polynomial of degree 20 over $\QQ$.
The primary verifier checks positive endpoint values and zero roots on the
complete interval by an exact rational Sturm sequence.  The independent
verifier reconstructs the numerator from 21 exact rational values and
converts it on $[0,1/50]$ to the Bernstein basis of degree 26; all 27
coefficients are strictly positive.  Both calculations also check
\eqref{eq:D-endpoint}.  In particular,
\[
 H(0)=L_{10}-\gamma_{10}=\frac{431}{10^8}>0.
\]
This excludes every marginal tight cut.

\section{Proper tight cuts}

Suppose now that $u,v\ge1$.  Define
\[
 \varepsilon_r=\sum_{i\in I^c}r_i-u,
 \qquad
 \varepsilon_c=\sum_{j\in J^c}c_j-v.
\]
For positive numbers $x_1,\dots,x_n$ of sum $n$ and product $X$, binary
Pinsker applied after grouping a subset and its complement gives
\begin{equation}\label{eq:subset}
 \left(\sum_{i\in S}x_i-|S|\right)^2
 \le-\frac n2\log X.
\end{equation}
Applying \eqref{eq:subset} to the row and column sums, then using
Cauchy--Schwarz and \eqref{eq:product-lower}, yields
\begin{equation}\label{eq:T}
 |\varepsilon_r|+|\varepsilon_c|
 \le\sqrt{-n\log(RC)}
 \le\sqrt{\frac{n\delta}{1-\delta}}=:T.
\end{equation}
The cut identity
\[
 s(A[I,J])=k-\varepsilon_r-\varepsilon_c+s(A[I^c,J^c])
\]
therefore implies
\begin{equation}\label{eq:lambda-proper}
 \lambda\ge1-\frac Tk.
\end{equation}
For $n=10$, the exact inequality $n\gamma_n<1-\gamma_n$ ensures $T<1\le k$.

Because $B$ is doubly stochastic,
\[
 s(B[I,J])=k+s(B[I^c,J^c])\ge k.
\]
Tightness and \eqref{eq:domination} give the reverse bound after
multiplication by $\lambda$, so equality holds and
\begin{equation}\label{eq:zero-block}
 B[I^c,J^c]=0.
\end{equation}
Thus $B$ lies on the staircase face with a prescribed $u\times v$ zero
block.  Hwang's theorem says that the barycenter minimizes the permanent on
this face.  Its permanent is
\begin{equation}\label{eq:nu}
 \per B\ge\nu_{n;u,v}
 :=\frac{\gamma_{n-u}\gamma_{n-v}}{\gamma_k},
 \qquad k=n-u-v,
\end{equation}
where $\gamma_0=1$.  Indeed, the barycenter has
$(n-v)!(n-u)!/k!$ supported permutations, all of the same weight; both
verifiers reconstruct this count exactly.

Using \eqref{eq:lambda-proper}, Bernoulli's inequality, and
$1-\delta\ge1-\gamma_n$ gives
\begin{align*}
 \nu_{n;u,v}\left(1-\frac Tk\right)^n-(\gamma_n-\delta)
 &\ge \nu_{n;u,v}-\gamma_n+\delta
 -\frac{n\nu_{n;u,v}}k
  \sqrt{\frac{n\delta}{1-\gamma_n}}.
\end{align*}
Completing the square in $\sqrt\delta$ gives the uniform lower bound
\begin{equation}\label{eq:eta}
 \eta_{n;u,v}
 :=\nu_{n;u,v}-\gamma_n
 -\frac{n^3\nu_{n;u,v}^2}{4k^2(1-\gamma_n)}.
\end{equation}
Both verifiers check $\eta_{10;u,v}>0$ for all
\[
 u,v\ge1,\qquad u+v\le9,
\]
namely all 36 proper cuts.  The smallest value occurs at $(u,v,k)=(1,1,8)$:
\[
 \nu_{10;1,1}=\frac{75161927680}{205891132094649},
\]
and
\[
 \eta_{10;1,1}
 =\frac{171312851843447795657879293024418789}
 {103456482841067610005769448332714895312500}>0.
\]
This contradicts $P=\gamma_n-\delta$ in every proper-cut case.

The whole, marginal, and proper cases exhaust every tight cut.  No boundary
global maximizer exists.  Hwang's positive-support theorem identifies the
sole remaining maximizer as $J_{10}/10$, proving
Theorem~\ref{thm:main}.

\section{Exact computational audit}

No floating-point value is used in a correctness decision.  The primary
standard-library verifier checks:
\begin{enumerate}
\item the bivariate equalization identity \eqref{eq:Ep};
\item the permanent polynomial \eqref{eq:Pn};
\item the exact lower bound \eqref{eq:L} by a rational Sturm sequence;
\item the degree-20 numerator of \eqref{eq:H}, denominator positivity,
endpoint positivity, and its exact zero root count;
\item the staircase barycenter count and all 36 values in \eqref{eq:eta}.
\end{enumerate}

The independent verifier does not import the primary verifier and shares
neither of its positivity procedures.  It uses a literal exact Ryser sum and
Lagrange interpolation for the two-zero polynomial, followed by Bernstein
positivity on 256 rational subintervals.  It reconstructs the marginal
numerator by exact value interpolation and proves it positive using a
degree-26 Bernstein expansion.  Its smallest certified Bernstein
coefficients are, respectively,
\[
 \frac{2098318974119444468161384333}
 {3961408125713216879677197516800000000}>0
\]
and
\[
 \frac{107851173950140635743821359}
 {39672851562500000000000000}>0.
\]
It separately checks the equalization identity, $D(1/50)>\gamma_{10}$, the
staircase counts, and all proper-cut inequalities.

The integration test runs both verifiers in ordinary and optimized Python
mode and requires identical output, confirming that no correctness decision
depends on an executable \texttt{assert} statement.  It also replaces
$L_{10}$ by a false inflated bound and requires both implementations to
reject the mutation.

\section{Scope and status}

This note establishes dimension 10.  Together with the other exact working
proofs in the accompanying repository, the currently covered dimensions are
\[
 4,5,10,11,12,13,14,15,16.
\]
Dimensions $6,7,8,9$ are not settled by these materials.  Pang's June 2026
preprint treats $n\ge17$ \cite{Pang2026}.  The present result should be
described as an exact computer-assisted working proof until the mathematical
reduction and both programs have undergone independent peer review.

\section*{Reproducibility}

From the \texttt{n10} directory, run
\begin{verbatim}
python3 -I verify_dittert_n10.py
python3 -I audit_dittert_n10_stdlib.py
python3 -I test_n10_package.py
\end{verbatim}
The verifier final lines must be
\begin{verbatim}
CERTIFIED
INDEPENDENT AUDIT CERTIFIED
\end{verbatim}
and the tests must report \texttt{OK}.

\begin{thebibliography}{99}

\bibitem{CheonWanless2012}
G.-S. Cheon and I. M. Wanless,
\emph{Some results towards the Dittert conjecture on permanents},
Linear Algebra Appl. 436 (2012), 791--801.
\href{https://doi.org/10.1016/j.laa.2010.08.041}{doi:10.1016/j.laa.2010.08.041}.

\bibitem{Egorychev1981}
G. P. Egorychev,
\emph{Solution of the van der Waerden problem for permanents},
Soviet Math. Dokl. 23 (1981), 619--622.

\bibitem{Falikman1981}
D. I. Falikman,
\emph{A proof of van der Waerden's conjecture on the permanent of a stochastic matrix},
Mat. Zametki 29 (1981), 931--938.

\bibitem{Hwang1985}
S.-G. Hwang,
\emph{Minimum permanent on faces of staircase type of the polytope of doubly stochastic matrices},
Linear Multilinear Algebra 18 (1985), 271--306.
\href{https://doi.org/10.1080/03081088508817694}{doi:10.1080/03081088508817694}.

\bibitem{Hwang1986}
S.-G. Hwang,
\emph{A note on a conjecture on permanents},
Linear Algebra Appl. 76 (1986), 31--44.
\href{https://doi.org/10.1016/0024-3795(86)90212-0}{doi:10.1016/0024-3795(86)90212-0}.

\bibitem{Li1990}
C.-K. Li,
\emph{On certain convex matrix sets},
Discrete Math. 79 (1990), 323--326.

\bibitem{Minc1984}
H. Minc,
\emph{Minimum permanents of doubly stochastic matrices with prescribed zero entries},
Linear Multilinear Algebra 15 (1984), 225--243.
\href{https://doi.org/10.1080/03081088408817592}{doi:10.1080/03081088408817592}.

\bibitem{Pang2026}
Z. Pang,
\emph{Proof of Dittert's conjecture for dimensions $n\ge17$},
arXiv:2606.01531 (2026).

\bibitem{PulaSongWanless2011}
K. Pula, S.-Z. Song, and I. M. Wanless,
\emph{Minimum permanents on two faces of the polytope of doubly stochastic matrices},
Linear Algebra Appl. 434 (2011), 232--238.
\href{https://doi.org/10.1016/j.laa.2010.08.022}{doi:10.1016/j.laa.2010.08.022}.

\end{thebibliography}

\end{document}
